

\section{Recursive ambiguity and persistent contracts}
Bayesian conditioning and a fresh worst-case choice at every date answer
different questions. Suppose a persistent binary type $\theta$ is either zero
or one, and tomorrow a fair binary signal $X$ unrelated to that type is observed.
The admitted laws are $Q_0$ and $Q_1$: under $Q_i$, $\theta=i$ with certainty
and $X$ remains fair. For the loss $Z=\ind\{X=\theta\}$, either law gives
expected loss $1/2$, so
\begin{equation}
\max_{Q\in\{Q_0,Q_1\}}\E_Q[Z]=\tfrac12,\qquad
\E_X\!\left[\max_{i\in\{0,1\}}\E_{Q_i}[Z\mid X]\right]=1.
\end{equation}
The second expression selects type zero after signal zero and type one after
signal one. It therefore evaluates a new law with $\theta=X$. That law is
absent even from the convex hull of the original two laws: every mixture
retains independence between the signal and type. No type has physically
changed along a realized path, yet dependence has changed across possible
paths. Preserving path support alone would miss the error.

Shapiro defines rectangularity through equality between static worst-case
evaluation and its nested conditional counterpart
\cite[Definition~2.1, Theorem~2.1 and proof, Section~5]{litcloseshapiro}.
His finite construction pastes conditional distributions across histories; it
can enlarge the family. A dynamic program for that enlarged family solves
a changed robustness question. Its convenience cannot justify claiming that
it solved the original persistent-type problem.

Terminology also requires care. Pichler and Shapiro call a product family a
rectangular setting, but explicitly show that its original worst-case
functional and its composite conditional functional can differ
\cite[Section~4.2, Example~4.7]{litclosepichlershapiro}. Computing a conditional
law-based risk measure can also differ from conditioning an unconditional
ambiguity set. Take six equiprobable states and an observed three-state
group with losses $0,1,2$. Conditional expected shortfall at level $1/2$
within that group is $5/3$. The unconditional six-state risk envelope permits
probabilities at most $1/3$ per state. It can assign $1/6$ to the loss-2 state,
zero to the other two group members, and $5/18$ to each state outside the
group. This is an admissible probability law; after conditioning on the group
it assigns probability one to loss 2. The conditional worst case over that
envelope is therefore 2. The difference is between two precise conditional
risk definitions, not an arithmetic error
\cite[Example~3.2 and equations~(3.1)--(3.5)]{litclosepichlershapiro}.

Policy consistency and equality of risk values are distinct as well. Xin and
Goldberg give moment-constrained inventory examples in which an initially
optimal policy ceases to be optimal after reoptimization, and another in which
the optimal policy sets agree although the two optimal values differ
\cite[Definitions~3.1--3.3, Observations~5--7, proofs in Sections~6.3--6.5]{litclosexingoldberg}.
The company model must specify whether management can commit to a policy,
what information each later decision uses, and which family of joint laws is
retained. For a finite reference problem, enumerate complete admissible laws
and observation-measurable policies before choosing an approximation. For
larger problems, a set of beliefs over persistent types or an explicit
precommitment formulation can retain those restrictions, but its recursive
sufficiency and computational cost need a separate derivation. Repeated
worst-case optimization remains useful as a stress test when the changed
family is stated explicitly.
